\chapter{ML-DSA AXI Wrapper Specification}

\section*{Purpose}
This chapter specifies the stable PoC AXI-Lite wrapper that presents the hardware signing engine to Linux software.

\section{Design Intent}
The wrapper provides a narrow, deterministic control and status surface between PS software and PL logic. It remains intentionally simple so that software, simulation, and future core refinement can converge on a stable contract.

\section{Current Implemented Structure}
The repository contains a stable AXI-Lite wrapper that instantiates a generic ML-DSA engine adapter. The wrapper owns the PS-visible register map and sequencing contract, while the adapter and downstream shim own engine-facing completion behavior.

\section{Primary Functions}
\begin{itemize}[leftmargin=*]
  \item Accept a 64-byte digest written by software.
  \item Expose control bits for start and clear-status behavior.
  \item Report idle, busy, done, and error status.
  \item Expose signature length and a read window for signature data.
  \item Remain agnostic to whether the adapter is operating in `STUB`, `CORE\_PLACEHOLDER`, or `MLDSA\_OSH` mode.
\end{itemize}

\section{Wrapper-to-Adapter Boundary}
The wrapper launches the adapter with a `start` event and a 64-byte digest, then waits for adapter completion. The adapter returns busy, done, error, signature length, and signature buffer information. The wrapper does not encode knowledge of ML-DSA-OSH internal state machines, stream ordering, or key segmentation.

\section{Current Behavioral Model}
The wrapper operates as a single-request engine for the current phase. Software shall not issue a new `start` command while the wrapper reports busy. A valid `start` clears stale visible results, launches the adapter, and later exposes a mode-dependent result once the adapter reports completion or error.

\section{Current Adapter Modes}
\begin{itemize}[leftmargin=*]
  \item \textbf{STUB:} fixed 128-byte signature, deterministic two-cycle completion latency, and signature content `STUBSIG || digest || zero padding`.
  \item \textbf{CORE\_PLACEHOLDER:} fixed 128-byte signature, deterministic four-cycle completion latency, and placeholder content `COREPH || digest || zero padding`.
  \item \textbf{MLDSA\_OSH:} real imported-core attachment path behind the adapter. When the real mixed-language core is available, the wrapper reports the actual ML-DSA-87 signature length. In the current local Verilog-only regression flow, this mode is tested through an honest fallback that reports engine error instead of claiming cryptographic success.
\end{itemize}

\section{Control Semantics}
\begin{itemize}[leftmargin=*]
  \item `start`: launches a new adapter transaction when the wrapper is not busy.
  \item `clear\_status`: clears done and error state, and clears the exposed signature buffer when the wrapper is idle.
\end{itemize}

\section{Error Handling}
The wrapper reports an error if software attempts to start a new transaction while busy or touches an unsupported register region. It also reports an engine error when the downstream adapter signals failure, including the current local `MLDSA\_OSH` fallback case used when the real imported core is not elaborated.

\section{Compatibility Note}
The register base offsets are unchanged from earlier phases. The signature-data window has been compatibly extended so the wrapper can surface the larger real ML-DSA-87 signature without introducing a new register-bank concept or changing software-visible control semantics.